\subsection{EXISTENCE AND UNIQUENESS THEOREM}

Let \(\lambda \in \mathfrak{h}^*\) . Since \(\mathfrak{b}_{+} = \mathfrak{h} \oplus \mathfrak{n}_{+}\) and since \(\mathfrak{n}_{+} = [\mathfrak{b}_{+}, \mathfrak{b}_{+}]\) , the map \(h + n \mapsto \lambda(h)\) (where \(h \in \mathfrak{h}\) , \(n \in \mathfrak{n}_{+}\) ) from \(\mathfrak{b}_{+}\) to \(k\) is a 1-dimensional representation of \(\mathfrak{b}_{+}\) . Denote by \(\mathrm{L}_{\lambda}\) the \(k\) -vector space \(k\) equipped with the \(\mathfrak{b}_{+}\) -module structure defined by this representation. Let \(\mathrm{U}(\mathfrak{g})\) , \(\mathrm{U}(\mathfrak{b}_{+})\) be the enveloping algebras of \(\mathfrak{g}\) , \(\mathfrak{b}_{+}\) , so that \(\mathrm{U}(\mathfrak{b}_{+})\) is a subalgebra of \(\mathrm{U}(\mathfrak{g})\) ; recall that \(\mathrm{U}(\mathfrak{g})\) is a free right \(\mathrm{U}(\mathfrak{b}_{+})\) -module (Chap. I, §2, no. 7, Cor. 5 of Th. 1). Put

\[
\mathrm{Z} (\lambda) = \mathrm{U} (\mathfrak {g}) \otimes_ {\mathrm{U} (\mathfrak {b} _ {+})} \mathrm{L} _ {\lambda}.\tag{2}
\]

Then \(Z(\lambda)\) is a left \(\mathfrak{g}\) -module. Denote by \(e\) the element \(1 \otimes 1\) of \(Z(\lambda)\) .

PROPOSITION 5. (i) The element \(e\) of \(Z(\lambda)\) is primitive of weight \(\lambda\) and generates the \(\mathfrak{g}\) -module \(Z(\lambda)\) .

\begin{enumerate}
\def\labelenumi{(\roman{enumi})}
\setcounter{enumi}{1}
\item
  Let \(\mathrm{Z}^{+}(\lambda) = \sum_{\lambda \neq \mu} \mathrm{Z}(\lambda)^{\mu}\) . Every submodule of \(\mathrm{Z}(\lambda)\) distinct from \(\mathrm{Z}(\lambda)\) is contained in \(\mathrm{Z}^{+}(\lambda)\) .
\item
  There exists a largest submodule \(F_{\lambda}\) of \(Z(\lambda)\) distinct from \(Z(\lambda)\) . The quotient module \(Z(\lambda)/F_{\lambda}\) is simple and has highest weight \(\lambda\) .
\end{enumerate}

It is clear that e generates the g-module \(Z(\lambda)\) . If \(x \in b_{+}\) ,

\[
x. e = (x. 1) \otimes 1 = (1. x) \otimes 1 = 1 \otimes x. 1 = \lambda (x) (1 \otimes 1) = \lambda (x) e,
\]

hence (i).

The \(\mathfrak{h}\) -module \(Z(\lambda)\) is semi-simple (Prop. 1). If \(G\) is a \(\mathfrak{g}\) -submodule of \(Z(\lambda)\) , then \(G = \sum_{\mu \in \mathfrak{h}^*}(G \cap Z(\lambda)^\mu)\) . The hypothesis \(G \cap Z(\lambda)^\lambda \neq 0\) implies that \(G = Z(\lambda)\) , since \(\dim Z(\lambda)^\lambda = 1\) and \(e\) generates the \(\mathfrak{g}\) -module \(Z(\lambda)\) . If \(G \neq Z(\lambda)\) , then \(G = \sum_{\mu \neq \lambda} G \cap Z(\lambda)^\mu \subset Z^+( \lambda)\) .

Let \(F_{\lambda}\) be the sum of the g-submodules of \(Z(\lambda)\) distinct from \(Z(\lambda)\) . By (ii), \(F_{\lambda} \subset Z^{+}(\lambda)\) . Hence \(F_{\lambda}\) is the largest submodule of \(Z(\lambda)\) distinct from \(Z(\lambda)\) . It is clear that \(Z(\lambda)/F_{\lambda}\) is simple and that the canonical image of e in \(Z(\lambda)/F_{\lambda}\) is primitive of weight \(\lambda\) .

In the remainder of this chapter, the g-module \(Z(\lambda)/F_{\lambda}\) of Prop. 5 will be denoted by \(E(\lambda)\) .

THEOREM 1. Let \(\lambda \in \mathfrak{h}^*\) . The \(\mathfrak{g}\) -module \(\mathrm{E}(\lambda)\) is simple and has highest weight \(\lambda\) . Every simple \(\mathfrak{g}\) -module of highest weight \(\lambda\) is isomorphic to \(\mathrm{E}(\lambda)\) .

The first assertion follows from Prop. 5 (iii). The second follows from Prop. 4.

PROPOSITION 6. Let V be a g-module, \(\lambda\) an element of \(h^{*}\) and v a primitive element of V of weight \(\lambda\) .

\begin{enumerate}
\def\labelenumi{(\roman{enumi})}
\item
  There exists a unique homomorphism of \(\mathfrak{g}\) -modules \(\psi: \mathrm{Z}(\lambda) \to \mathrm{V}\) such that \(\psi(e) = v\) .
\item
  Assume that v generates V. Then \(\psi\) is surjective. Moreover, \(\psi\) is bijective if and only if, for every non-zero element u of \(\mathrm{U}(\mathfrak{n}_{-})\) , \(u_{V}\) is injective.
\item
  The map \(u \mapsto u \otimes 1\) from \(\mathrm{U}(\mathfrak{n}_{-})\) to \(\mathrm{Z}(\lambda)\) is bijective.
\end{enumerate}

Let K be the kernel of the representation of \(\mathrm{U}(\mathfrak{b}_{+})\) on \(L_{\lambda}\) ; it is of codimension 1 in \(\mathrm{U}(\mathfrak{b}_{+})\) . Let \(\mathrm{J} = \mathrm{U}(\mathfrak{g})\mathrm{K}\) be the left ideal of \(\mathrm{U}(\mathfrak{g})\) generated by K; then \(L_{\lambda}\) can be identified with \(\mathrm{U}(\mathfrak{b}_{+})/\mathrm{K}\) as a left \(\mathrm{U}(\mathfrak{b}_{+})\) -module, and \(\mathrm{Z}(\lambda)\) can be identified with \(\mathrm{U}(\mathfrak{g})/\mathrm{J}\) as a left \(\mathrm{U}(\mathfrak{g})\) -module. We have K.v = 0, so J.v = 0, which proves (i).

Now assume that \(v\) generates V. It is clear that \(\psi\) is surjective.

By the Poincaré-Birkhoff-Witt theorem (Chap. I, §2, no. 7, Cor. 6 of Th. 1), a basis of \(\mathrm{U}(\mathfrak{n}_{-})\) over \(k\) is also a basis of \(\mathrm{U}(\mathfrak{g})\) as a right \(\mathrm{U}(\mathfrak{b}_{+})\) -module. Hence the map \(\varphi : u \mapsto u \otimes 1\) from \(\mathrm{U}(\mathfrak{n}_{-})\) to \(\mathrm{U}(\mathfrak{g}) \otimes_{\mathrm{U}(\mathfrak{b}_{+})} \mathrm{L}_{\lambda}\) is bijective. Let \(u \in \mathrm{U}(\mathfrak{n}_{-})\) . Then \(\varphi^{-1} \circ u_{\mathrm{Z}(\lambda)} \circ \varphi\) is left multiplication by \(u\) on \(\mathrm{U}(\mathfrak{n}_{-})\) . In view of Chap. I, §2, no. 7, Cor. 7 of Th. 1, \(u_{\mathrm{Z}(\lambda)}\) is injective if \(u \neq 0\) . Consequently, if \(\psi\) is bijective, then \(u_{\mathrm{V}}\) is injective for non-zero \(u\) in \(\mathrm{U}(\mathfrak{n}_{-})\) .

Assume that \(\psi\) is not injective. There exists \(u \in \mathrm{U}(\mathfrak{n}_{-})\) such that \(u \neq 0\) and \(\psi(\varphi(u)) = 0\) . Then

\[
u _ {\mathrm{V}}. v = u _ {\mathrm{V}}. \psi (1 \otimes 1) = \psi (u \otimes 1) = \psi (\varphi (u)) = 0.
\]

COROLLARY 1. Let \(\lambda \in \mathfrak{h}^*\) and \(\alpha \in B\) be such that \(\lambda(H_{\alpha}) + 1 \in \mathbf{N}\) . Then \(Z(-\alpha + s_{\alpha}\lambda)\) is isomorphic to a \(\mathfrak{g}\) -submodule of \(Z(\lambda)\) .

Put \(m = \lambda(H_{\alpha})\) . Let \(x = X_{-\alpha}^{m+1}.e \in Z(\lambda)\) , and let V be the submodule of \(Z(\lambda)\) generated by x. Then \(x \neq 0\) (Prop. 6). On the other hand, \(x \in Z(\lambda)^{\lambda-(m+1)\alpha}\) . For \(\beta \in B\) and \(\beta \neq \alpha\) , \([\mathfrak{g}^{-\alpha}, \mathfrak{g}^{\beta}] = 0\) and \(g^{\beta}.e = 0\) , so \(g^{\beta}.x = 0\) . Finally, since \([X_{\alpha}, X_{-\alpha}] = -H_{\alpha}\) , we have

\[
[ X _ {\alpha}, X _ {- \alpha} ^ {m + 1} ] = (m + 1) X _ {- \alpha} ^ {m} (- H _ {\alpha} + m)
\]

(§1, no. 1, Lemma 1 (ii)), so

\[
X _ {\alpha}. x = X _ {\alpha} X _ {- \alpha} ^ {m + 1}. e = [ X _ {\alpha}, X _ {- \alpha} ^ {m + 1} ]. e = (m + 1) X _ {- \alpha} ^ {m} (m e - \lambda (H _ {\alpha}) e) = 0.
\]

Thus, \(x\) is primitive of weight \(\lambda - (m + 1)\alpha\) . In view of Prop. 6, the \(\mathfrak{g}\) -module \(V\) is isomorphic to \(Z(-\alpha + \lambda - m\alpha) = Z(-\alpha + s_{\alpha}\lambda)\) .

COROLLARY 2. Let \(\rho = \frac{1}{2}\sum_{\alpha \in \mathrm{R}_{+}}\alpha\) , and \(\lambda, \mu \in \mathfrak{h}^*\) . Assume that \(\lambda + \rho\) is a dominant weight in \(\mathrm{R}\) , and that there exists \(w \in \mathrm{W}\) with \(\mu + \rho = w(\lambda + \rho)\) . Then \(\mathrm{Z}(\mu)\) is isomorphic to a submodule of \(\mathrm{Z}(\lambda)\) .

The assertion is clear when w = 1. Assume that it is established whenever w is of length \textless{} q. If w is of length q, there exists \(\alpha \in B\) such that \(w = s_{\alpha}w'^{-1}\) , with \(l(w') = q - 1\) . We have \(w'(\alpha) \in \mathrm{R}_{+}\) (Chap. VI, §1, no. 6, Cor. 2 of Prop. 17), and hence \(w'^{-1}(\lambda + \rho)(H_{\alpha}) = (\lambda + \rho)(H_{w'\alpha})\) is an integer \(\geq 0\) . Put

\[
\mu^ {\prime} = w ^ {- 1} (\lambda + \rho) - \rho .
\]

By the induction hypothesis, \(Z(\mu')\) is isomorphic to a submodule of \(Z(\lambda)\) . On the other hand, by Chap. VI, §1, no. 10, Prop. 29 (ii),

\[
- \alpha + s _ {\alpha} \mu^ {\prime} = - \alpha + s _ {\alpha} w ^ {- 1} (\lambda + \rho) - s _ {\alpha} \rho = w (\lambda + \rho) - \rho = \mu .
\]

Moreover, \(\rho(H_{\alpha}) = 1\) (Chap. VI, §1, Prop. 29 (iii)), so \(\mu'(H_{\alpha}) + 1 \in \mathbf{N}\) . Cor. 1 now implies that \(Z(\mu)\) is isomorphic to a submodule of \(Z(\mu')\) , and hence also to a submodule of \(Z(\lambda)\) .

